% Declaración exigida por el Instructivo N.º 17/2026 (VAA, Universidad de Chile): páginas preliminares.
% El detalle de prompts y metodología va en el Anexo chap:anexoIA (capitulos/anexo_ia.tex).
\chapter*{Declaración de uso de inteligencia artificial}

En la elaboración de esta tesis se utilizaron herramientas de inteligencia artificial generativa. La principal fue Claude (Anthropic), en sus versiones Sonnet~5, Opus~5 y Opus~5.5, usada entre agosto y octubre de 2026 a través de Claude Code, un asistente que lee y edita los archivos del proyecto y ejecuta código en el computador del autor. De forma secundaria se usó el asistente de inteligencia artificial del editor \LaTeX{} Prism para ajustar diapositivas de la presentación.
% TODO (autor): confirmar el proveedor y el uso del asistente de Prism, y agregar cualquier otra
% herramienta usada fuera de este repositorio (antes de agosto de 2026 o en otros equipos).

Estas herramientas se emplearon en seis etapas: (i)~\emph{generación de código}, para la extracción de datos de Bloomberg, la construcción del panel, la estimación de $JLoss$ y del $GaR$, las baterías de regresiones y las pruebas de robustez; (ii)~\emph{análisis de datos}, para ejecutar esos programas y producir tablas y figuras; (iii)~\emph{revisión de literatura}, para sintetizar artículos y preparar fichas de lectura; (iv)~\emph{redacción y edición} de la introducción, el marco teórico, la metodología, los resultados, la discusión y las conclusiones, incluidas la corrección editorial, la consistencia de las cifras y el formato \LaTeX; (v)~\emph{revisión crítica simulada}, con informes de tipo árbitro que propusieron pruebas adicionales; y (vi)~la \emph{preparación de la defensa}, con láminas y guiones. Como la inteligencia artificial intervino en componentes fundamentales de la tesis, los prompts representativos y la metodología de integración y verificación se describen en el Anexo~\ref{chap:anexoIA}.

La pregunta de investigación, las hipótesis, las decisiones sobre la muestra y la especificación, y la interpretación de los resultados son del autor, discutidas con su profesor guía. Todo el contenido fue revisado y validado por el autor, quien asume plena responsabilidad sobre el trabajo presentado.
